\documentclass[]{llncs}
\usepackage{xcolor}
\usepackage{pdfpages}
\usepackage{longtable}
\newcommand{\tocTitle}[2]{\contentsline{title}{#1}{#2}{}}
\newcommand{\tocAuthors}[1]{{\raggedright \leftskip 15pt \rightskip 2.55em\itshape #1\endgraf}}
\newcommand{\meta}[1]{{\color{blue}#1}}
\newcommand{\tocSection}[1]{\subsection*{#1}}
\pagestyle{plain}
\pagenumbering{roman}
\setcounter{page}{5}

\sloppypar

\begin{document}


\chapter*{Preface}
\markboth{Preface}{Preface}

This volume contains the proceedings of the 28th International Conference on Coordination Models and Languages (COORDINATION 2026), held on June 8–12, 2026, in Urbino, Italy hosted by the
University of Urbino Carlo Bo., as part of the 21st International Federated
Conference on Distributed Computing Techniques (DisCoTec 2026).


Modern information systems rely increasingly on combining concurrent, distributed,
mobile, adaptive, reconfigurable, and heterogeneous components. New models, architectures, languages, and verification techniques are necessary to cope with the complexity induced by the demands of today’s software development. Coordination languages
have emerged as a successful approach, in that they provide abstractions that cleanly
separate behavior from communication, thereby increasing modularity, simplifying
reasoning, and ultimately enhancing software development. Building on the success of	
the previous editions, COORDINATION 2026 provides a well-established forum for the
growing community of researchers interested in models, languages, architectures, and
implementation techniques for coordination.

COORDINATION 2026 solicited contributions in three different categories: (1) regular papers describing thorough and complete research results and experience reports;
(2) survey papers describing important results and success stories related to
the topics of COORDINATION; and (3) tool papers describing technological artifacts
in the scope of the research topics of COORDINATION, and requiring the submission
of the associated artifact via a separate procedure.

%There were 41 submissions. Each submission was reviewed by at least 0 people, and on the average 2.6 people, program committee members. The committee decided to accept 13 papers. The program also includes 0 talks.

There were 28 paper submissions sent for peer-review distributed over the different categories: 23 regular
papers, and 5 tool papers. The selection of the papers
was entrusted to the Programme Committee (PC), consisting of 32 members from 15
different countries (with gender balance F/M = 12/20). The selection of the papers was
done electronically, in two phases. In the first phase, which lasted three weeks, each
submission was single-blind reviewed by at least three PC members, in some cases with
the help of external reviewers. During the second phase, which lasted four days, the papers were thoroughly discussed. The decision to accept or reject a paper
was based not only on the reviews and scores but also on these in-depth discussions. In
the end, 13 papers were selected to be presented at the conference, for an acceptance ratio of 47\%: 11 regular papers, and 2 tool papers. 


The authors of the submitted tool papers were requested to participate in the EAPLS
artifact badging, which instead was optional for all the other paper categories. There
were 13 artifact submissions. The 9 members of the Artifact Evaluation Committee
(AEC), chaired by Gianluca Aguzzi, awarded the Available and Reusable Badge to 3
artifacts, and the Available and Functional Badge to 4 artifacts.

%You can write the history of your conference here.

We organized the accepted papers into five groups according to the topics they address:

\paragraph{Theoretical foundations and models for distributed systems.} 
This group, containing three papers, focuses on formal reasoning of distributed system properties. 
\begin{itemize}
\item The paper ``Timed Scenario Expressions and Realisability''  addresses the problem of realisability for the distributed system modeled by sets of sequential timed scenarios.  
\item The paper ``HistMSO: a Logic for Reasoning on Consistency Models with MONA'' reasons on the consistency of replicated data systems modeled by sequences of data operations. 
\item The paper ``Motif Refinement for the Hierarchical Control of Structured CPSs" introduces a refinement-based modeling and development approach for cyber-physical systems which preserves control stability.  
\end{itemize}

\paragraph{Distributed algorithms for collective adaptive systems.} 
This group contains two papers which reformulate collective distributed algorithms within the aggregate programming framework with specific guarantees. 
\begin{itemize}
\item The paper ``Aggregate Indoor Localisation" provides two existing established localisation algorithms to guarantees  adaptability and resilience while improving their performance. 
\item The paper ``A Self-Stabilizing Min-Max Consensus via Path-loop Detection" implements the min-max consensus algorithm as a
	reusable library for the Collektive DSL, which guarantees the self-stabilization property. 
\end{itemize}

\paragraph{Coordination languages, frameworks, and patterns.}
This group includes contributions centered on the design and implementation of coordinated systems through programming languages and middleware support.
\begin{itemize}
\item 	The paper ``ScalaTropy: Multiparty Coordination with Monadic Communication Primitives'' provides a design and Scala-internal DSL that combines choreographic and multi-tier programming to support multiple coordination patterns.
\item The paper ``Bach4Popper: Towards Federated Inductive Logic Programming using Coordination'' proposes a distributed framework and architecture for federated inductive logic programming, to support learning while ensuring explainability and privacy.
\end{itemize}


\paragraph{Formal verification in coordination languages.}
%
This group focuses on the static verification of two classes of coordination languages. 
\begin{itemize}
\item The paper ``Proof of Delivery: Mechanized Mailbox Types" introduces a machine-checked foundation for the programming languages, specifically Pat, that incorporate mailbox type system to reason about the contents of mailboxes. 
\item The paper ``Deductive Verification of Legal Contracts" enables deductive verification of contracts specified by Stipula through their translation into Java programs annotated with Java 
	Modeling Language specifications preserving the dynamic
	features of the original programs – most notably the scheduling of events at future times.
\end{itemize}

\paragraph{Simulation tools and methodologies for collective adaptive systems.} Tis group focuses on the simulation of systems consisting of a multitude of devices or agents. 
\begin{itemize}
\item The paper   ``Simulation and Analysis of Indoor-Air-Quality Measuring Devices with YODA'' provides an experience report on the design of an Internet of Things system using the agent-based language YODA and the Sibilla simulator for collective systems. \item Finally, the tool paper ``High-Fidelity Simulation of Aggregate Computing Systems with Collektivity'' presents an integration of Collektive and Unity for the simulation of aggregate computing systems.
\end{itemize}


As Programme Chairs, we actively contributed to the selection of the three keynote
speakers of DisCoTec 2026:\begin{itemize}
	\item Prof. Mehdi Dastani (Utrecht University, The Netherlands);
	\item Paolo Romano (NESC-ID, Portugal);
	\item Nathalie Bertrand (Inria, France).
\end{itemize}
We are most grateful to Prof. Dastani for accepting our invitation as the  COORDINATION-related keynote speaker. This volume includes the abstract of his keynote
talk: “Efficient Communication and Coordination in Multiagent Reinforcement Learning”.


As is traditional in DisCoTec, a joint session with the best papers from each main conference was organized. The best paper of COORDINATION 2026 was “Runtime adaptation as a programming pattern in service composition” authored by Carlos Gustavo Lopez Pombo, Pablo Montepagano and Emilio Tuosto. 

%Write any acknowledgements ...
We are grateful to all the persons involved in COORDINATION 2026. In particular,
we thank the authors for their submissions, the attendees of the conference for their
participation, the PC members and external reviewers for their work in reviewing sub-
missions and participating in the discussions, the AEC chair Gianluca Aguzzi and
the AEC members for their effort in the evaluation of the artifacts, and the Steering Committee,
chaired by Michele Loreti, for their guidance and support. We are also grateful to the
Organizing Committee, chaired by Claudio Antares Mezzina, for their excellent job. We also thank
the providers of the EasyChair Conference Management System, which was of great
help in the submission and reviewing process and in the preparation of the proceedings.
We would also like to acknowledge the prompt and professional support from Springer,
which published these proceedings in printed and electronic volumes as part of their
LNCS book series.


%Do not forget to mention your sponsors!
%We appreciate some words about the use of EasyChair but it is not necessary.



~\bigskip

\noindent
\begin{minipage}[t]{.4\textwidth}
April, 2026\\
%Cesena
\end{minipage}%
\hfill
\begin{minipage}[t]{.4\textwidth}\flushright
Roberto Casadei\\
Fatemeh Ghassemi
\end{minipage}


%~\bigskip
%
%\noindent
%\begin{minipage}[t]{.4\textwidth}
%April 19, 2026\\
%Cesena
%\end{minipage}%
%\hfill
%\begin{minipage}[t]{.4\textwidth}\flushright
%Roberto Casadei\\
%Fatemeh Ghassemi
%\end{minipage}

\clearpage


\chapter*{Organization}
\markboth{Organization}{Organization}
\setlength{\tabcolsep}{0pt}     % remove spacing between columns in the following long tables
\setlength{\LTpre}{0pt}         % remove extra vertical space before the long tables

\section*{DisCoTec General Chair}
\begin{longtable}{p{0.35\textwidth}p{0.65\textwidth}}
Claudio Antares Mezzina          & University of Urbino, Italy
\end{longtable}

\section*{DisCoTec Organising Committee}
\begin{longtable}{p{0.55\textwidth}p{0.45\textwidth}}
Sara Montagna (Gender-equality) & University of Urbino, Italy\\
Pierluigi Graziani (Local Organizer) & University of Urbino, Italy\\
Andrea Esposito (Satellite Events Chair) & University of Urbino, Italy\\
Claudio Antares Mezzina (General Chair) & University of Urbino, Italy
\\
Martin Vassor (Publicity Chair) & Universit\'e de Lorraine, France
\end{longtable}

\section*{DisCoTec Steering Committee}
\begin{longtable}{p{0.35\textwidth}p{0.65\textwidth}}
Simon Bliudze & Inria Centre at the University of Lille, France\\
Rocco De Nicola & IMT School for Advanced Studies Lucca, Italy\\
Adrian Francalanza (Chair) & University of Malta, Malta\\
Ivan Lanese & University of Bologna/Inria, Italy\\
Alberto Lluch Lafuente & Technical University of Denmark, Denmark\\
Michele Loreti & University of Camerino, Italy\\
Lu\'is Veiga & INESC-ID, Universidade de Lisboa, Portugal\\
Gianluigi Zavattaro & University of Bologna, Italy\\
Carla Ferreira & NOVA University of Lisbon, Portugal\\
Jorge A. P\'erez & University of Groningen, The Netherlands
\end{longtable}

\clearpage


\section*{COORDINATION Steering Committee}
\begin{longtable}{p{0.40\textwidth}p{0.60\textwidth}}
Gul Agha & University of Illinois at Urbana Champaign, USA\\
Farhad Arbab & CWI and Leiden University, The Netherlands\\
Simon Bliudze & INRIA Lille, France\\
Laura Bocchi & University of Kent, UK\\
Ilaria Castellani & INRIA, France\\
Ferruccio Damiani & University of Turin, Italy\\
Ornela Dardha & University of Glasgow, UK\\
Wolfgang De Meuter & Vrije Universiteit Brussel, Belgium\\
Rocco De Nicola & IMT School for Advanced Studies Lucca, Italy\\
Giovanna di Marzo Serugendo & Universit\'e de Gen\`eve, Switzerland\\
Tom Holvoet & KU Leuven, Belgium\\
Jean-Marie Jacquet & University of Namur, Belgium\\
Sung-Shik Jongmans & Open University of the Netherlands, Netherlands\\
Christine Julien & University of Texas at Austin, USA\\
Eva K\"uhn & Vienna University of Technology, Austria\\
Alberto Lluch Lafuente & Technical University of Denmark, Denmark\\
Ant\'onia Lopes & University of Lisbon, Portugal\\
Michele Loreti (Chair) & Universit\`a di Camerino, Italy\\
Mieke Massink & CNR-ISTI, Pisa, Italy\\
Jos\'e Proen\c{c}a & University of Porto, Portugal\\
Marjan Sirjani & M\"alardalen University, Sweden\\
Carolyn Talcott & SRI International, USA\\
Maurice ter Beek & CNR-ISTI, Italy\\
Francesco Tiezzi & Universit\`a di Firenze, Italy\\
Emilio Tuosto & Gran Sasso Science Institute, Italy\\
Vasco T. Vasconcelos & University of Lisbon, Portugal\\
Mirko Viroli & Universit\`a di Bologna, Italy\\
Gianluigi Zavattaro & Universit\`a di Bologna, Italy
\end{longtable}

\section*{COORDINATION Program Committee Chairs}

\noindent
\begin{longtable}{p{0.45\textwidth}p{0.55\textwidth}}
Roberto Casadei & University of Bologna, Italy\\

Fatemeh Ghassemi & University of Tehran, Itan\\
\end{longtable}


\section*{COORDINATION Program Committee}
\noindent
\begin{longtable}{p{0.45\textwidth}p{0.55\textwidth}}
S. Akshay & IIT Bombay, India\\

Duncan Paul Attard & University of Malta, Malta\\

Giorgio Audrito & University of Turin, Italy\\

Simon Bliudze & INRIA, France\\

Roberto Casadei & University of Bologna, Italy\\

Valentina Castiglioni & Eindhoven University of Technology, Netherlands\\

Mariangiola Dezani-Ciancaglini & University of Turing, Italy\\

Cinzia Di Giusto & Universit\'e C\^ote d'Azur, France\\

Francisco Ferreira & Imperial College, United Kingdom\\

Fatemeh Ghassemi & University of Tehran, Iran\\

Silvia Ghilezan & University of Novi Sad, Serbia\\

Susanne Graf & Universite Joseph Fourier / CNRS / VERIMAG, France\\

Ludovic Henrio & CNRS, France\\

Thomas Hildebrandt & University of Copenhagen, Denmark\\

Ping Hou & University of Oxford, United Kingdom\\

Eva K\"uhn & Vienna University of Technology, Austria\\

Ivan Lanese & University of Bologna, Italy\\

Carlos Gustavo Lopez Pombo & National University of Río Negro, Argentina\\

Michele Loreti & University of Camerino, Italy\\


Stefano Mariani & University of Modena and Reggio Emilia, Italy\\

Hernan Melgratti & Departamento Computaci\'on, University of Buenos Aires, Argentina\\

%Claudio Antares Mezzina & University of Urbino, Italy\\

Fabrizio Montesi & University of Southern Denmark, Denmark\\

J. Garrett Morris & The University of Kansas, United States of America\\

Rumyana Neykova & Brunel University, United Kingdom\\

Jos\'e Proen\c{c}a & University of Porto, Portugal\\

Violet Ka I Pun & Western Norway University of Applied Sciences, Norway\\

Ant\'onio Ravara & NOVA University Lisbon, Portugal\\

Marjan Sirjani & Malardalen University, Sweden\\

Carolyn Talcott & SRI  International, United States of America\\

Maurice ter Beek & ISTI-CNR, Italy\\

Francesco Tiezzi & University of Florence, Italy\\

Nils Timm & University of Pretoria, South Africa\\

Emilio Tuosto & Gran Sasso Science Institute, Italy\\

Frank Valencia & École Polytechnique de Paris, France\\

\end{longtable}

\clearpage

\section*{COORDINATION Artifact Evaluation Committee Chair}

\noindent
\begin{longtable}{p{0.35\textwidth}p{0.65\textwidth}}

Gianluca Aguzzi & University of Bologna, Italy

\end{longtable}

\section*{COORDINATION Artifact Evaluation Committee}

\noindent
\begin{longtable}{p{0.35\textwidth}p{0.65\textwidth}}

%Gianluca Aguzzi & University of Bologna, Italy\\

Davide Domini & University of Bologna, Italy\\

Nicolas Farabegoli & University of Bologna, Italy\\

Matteo Magnini & University of Luxembourg, Luxembourg\\

M\'ario Pereira & NOVA University Lisbon, Portugal\\


Marco Quadrini & University of Camerino, Italy\\

Corentin Reuther & University of Namur, Belgium\\

Lorenzo Rossi & University of Camerino, Italy\\

Gerard Tabone & University of Malta, Malta\\

Laura Voinea & University of Glasgow, United Kingdom\\

\end{longtable}


\section*{COORDINATION Additional Reviewers}
\noindent
\begin{longtable}{p{0.45\textwidth}p{0.55\textwidth}}
%\textbf{B} \\*
Davide Basile & ISTI-CNR, Italy\\
Nicola Del Giudice & University of Camerino, Italy\\
Masoud Ebrahimi & M\"alardalen University, Sweden\\
Gabriele Genovese & Universit\'e C\^ote d'Azur, France\\
Ramchandra Phawade & IIT Dharwad, India\\
Jos\'e Ignacio Requeno Jarabo & Complutense University of Madrid, Spain\\
Volker Stolz & H\o gskulen p\aa\ Vestlandet, Norway
\\
\end{longtable}




\clearpage

\chapter*{Invited contributions}


\includepdf{invited-talk.pdf}


\clearpage

\chapter*{Table of Contents}
\tocSection{Theoretical foundations and models for distributed systems}
  \tocTitle{Timed Scenario Expressions and Realisability}{1}
  \tocAuthors{Neda Saeedloei and Feliks Kluzniak}
  \tocTitle{HistMSO: a Logic for Reasoning about Consistency Models with MONA}{22}
  \tocAuthors{Isabelle Coget and \'Etienne Lozes}
  \tocTitle{Motif Refinement for the Hierarchical Control of Structured CPSs}{44}
  \tocAuthors{Simon Bliudze, Sophie Cerf and Olga Kouchnarenko}
\tocSection{Distributed algorithms for collective adaptive systems}
  \tocTitle{Aggregate Indoor Localisation}{67}
  \tocAuthors{Giorgio Audrito, Leonardo Bertolino, Ferruccio Damiani and Gianluca Torta}
  \tocTitle{A Self-Stabilizing Min-Max Consensus via Path-loop Detection}{87}
  \tocAuthors{Angela Cortecchia, Danilo Pianini and Mirko Viroli}
\tocSection{Coordination languages, frameworks, and patterns}
  \tocTitle{ScalaTropy: Multiparty Coordination with Monadic Communication Primitives}{110}
  \tocAuthors{Nicolas Farabegoli, Luca Tassinari, Gianluca Aguzzi and Mirko Viroli}
  \tocTitle{Bach4Popper: Towards Federated Inductive Logic Programming using Coordination}{131}
  \tocAuthors{Yasmine Akaichi, Manel Barkallah, Jean-Marie Jacquet, Isabelle Linden and Wim Vanhoof}
  \tocTitle{Runtime adaptation as a programming pattern in service composition}{153}
  \tocAuthors{Carlos Gustavo Lopez Pombo, Pablo Montepagano and Emilio Tuosto}
  \tocTitle{Phyelds: A Pythonic Framework for Aggregate Computing}{171}
  \tocAuthors{Gianluca Aguzzi, Davide Domini, Nicolas Farabegoli and Mirko Viroli}
\tocSection{Formal verification in coordination languages}
  \tocTitle{Proof of Delivery: Mechanized Mailbox Types}{190}
  \tocAuthors{Edgard Schiebelbein, Annette Bieniusa and Simon Fowler}
  \tocTitle{Deductive Verification of Legal Contracts}{211}
  \tocAuthors{Reiner H\"ahnle and Cosimo Laneve}
\tocSection{Simulation tools and methodologies for collective adaptive systems}
  \tocTitle{Simulation and Analysis of Indoor-Air-Quality Measuring Devices with YODA}{233}
  \tocAuthors{Riccardo Petracci, Nicola Del Giudice, Diletta Romana Cacciagrano and Michele Loreti}
  \tocTitle{High-Fidelity Simulation of Aggregate Computing Systems with Collektivity}{255}
  \tocAuthors{Filippo Gurioli, Martina Baiardi, Angela Cortecchia and Danilo Pianini}


\end{document}
